\begin{examplebox}
\paragraph{例1：求出候选点，再验两侧符号}
求 $f(x)=x^3-3x+1$ 的极值。由
\[
f'(x)=3(x-1)(x+1)
\]
知导数在 $(-\infty,-1),(-1,1),(1,+\infty)$ 上依次为 $+,-,+$。
所以 $x=-1$ 是极大值点，极大值 $f(-1)=3$；$x=1$ 是极小值点，极小值 $f(1)=-1$。
回验可直接展开：
\[
f(-1+h)-3=h^2(h-3)<0\quad(0<|h|<1),
\]
\[
f(1+h)+1=h^2(h+3)>0\quad(0<|h|<1).
\]

\paragraph{例2：不求导也能判断尖点极小值}
$g(x)=|x-2|+1$ 满足 $g(x)\ge1$，等号仅在 $x=2$ 成立，故 $2$ 是严格极小值点，极小值为 $1$，也是全局最小值。
回验：$g(2+h)-g(2)=|h|>0$（$h\ne0$）；左右导数分别为 $-1,1$，该点不可导。
\end{examplebox}
